\section{熊野寮生活あれこれ}
\rightline{（文責:まめめっち）}

\subsection{無関係な序文}
私は女子寮生です。学部時代を寮で過ごし、大学院に上がった今もお世話になっています。
ちなみに、この自己紹介は本文には全く関係ありません。

京都大学および熊野寮は全構成員のうち女性の占める割合が２割程度ととても少ないです\footnote{2025年の女子学生の割合は学部が22.8％で大学院が29.0％で概ね横ばいです。(京都大学.(2025).京都大学アカウンタビリティレポート2025.URL:https://www.kyoto-u.ac.jp/ja/about/data)}。そのため、寮に女性が住んでいるというと驚かれる人も多いです。

性別を明らかにしない場では熊野寮に住んでいるというだけで男性だと思われる。逆に、女であることが明らかな場では寮生に混じった寮外生だと思われる。なぜでしょう。

熊野寮生は全員男性に違いない（⇔女性なら熊野寮生のはずがない）という無意識のバイアスに、ジェンダーロールであることがなあと時々頭を抱えます。熊野寮には多くはないけど女子寮生が住んでいて、ごく当たり前にいろいろ考えながら生活をしているということを、少しでも感じてくれたらうれしいです。
あえてノイズを提供いたしました。

\subsection{本文}
以下では、熊野寮での生活ってどんな感じ？という疑問に対して私の答えをありのままに書いていこうと思います。良いことも悪いことも書きますので、もし熊野寮に夢をみている人がいたら、すこし壊してしまうかも。

皆さんご存じの通り、熊野寮は共同生活の場です。この「共同生活」というのは人によってかなり受け止め方が違うと思います。家族と同じ部屋で寝ているとか大人数の合宿をしたことがあるとか、そういったものとはかなり毛色の違う共同生活になると思います。熊野寮では、初めましてこんにちはの400人あまりと様々なモノを共有するのです。物品も、時間も、場所も。しかもかなり近い距離で。
\vskip.5\baselineskip
物品の共有についてはメリットも大きいです。なぜかというと、入居にあたって用意するものが少なく済み、初期費用がかなり抑えられるからです。私の先輩には布団すら持たずに入寮した人がいるそうです。私が入寮した時も、布団とスーツケース1つで入寮し、あとはドラッグストアで日用品を買うだけで生活できました。入寮初日は買い出しが間に合わなかったので、先輩寮生からお風呂セットを借りました。こういう互助が生まれるのも寮ならではなのかな？ありがたいことです。各ブロックには共用の調理器具があり、家電も概ね揃っています。自分が生活のために買うのは本当に消耗品くらいです。

一方で、他人とモノを共有することによる距離の近さがストレスになることもあります。例えば、私がよく利用する炊事場\footnote{炊事場とは、各階に2つ(A,B棟)か1つ（C棟）存在する、寮生共用の台所のこと。そのフロアやブロックの寮生が主な使用者となる}高確率で誰かの使用済み食器や開封済みの食品が放置されています。排水溝はごみで詰まり、スポンジは油で真っ黒ということも日常茶飯事。寮ではそういう生活上のトラブルを話し合いで解決することになっています\footnote{熊野寮内における徹底討論の原則}が、炊事場のように多くの人が使う場所は、その利用者が必ずしもお話し合いの場に来てくれるわけでもありません。私個人としてはこういうのに結構消耗してしまうのですが、多くの人は仕方がないやと心の距離をとって対処しているみたいです。
\vskip.5\baselineskip
次に、場所という観点で言うと、熊野寮でプライベートの確保には期待しない方がいいです。正直なところ、寮の中で一人になれる場所はベッドの上とトイレやシャワーの個室くらいだと思います。それらの場所でも、同部屋の人の出入りや順番待ちをしている他の寮生の存在が気になります。月々4300円はまあこんなものだろうと思います。ただ、一人になりたい時は、行き場に困ります。そういう時は、外に出る元気があったら、図書館や鴨川や近くの喫茶店\footnote{蛇足ですが、筆者のお気に入りカフェはcafe\&dining markです。寮から徒歩5分。日替わりランチの大盛ごはんで有名}で深呼吸したりしています。

寮生個人に割り当てられているものは、部屋によりますが机１つ、棚１つ、クローゼット１つ、ベッド１つ。物が少ないを自認している私でもちょっと収納には苦労しています。荷物が多い寮生は部屋の中で自分の場所を確保・整備するのに苦労している印象。もし実家に物をおいておけるなら、その方が楽かもしれないです。
\vskip.5\baselineskip
寮は広い建物ですが、それなりの人が住んでいるからやっぱりちょっと手狭かなと思います。でも、空間や時間を共有することで、いろいろな楽しいことができることも事実。他人と一緒になにかをやるハードルが低いのが熊野寮だと思います。「やってみたい」と言うと「じゃあやろう」という声が返ってくるし、他人の「やってみたい」に「わたしも」と乗れる雰囲気があります。

クリエイティブみたいな言葉とは無縁な私だったけれども、運動会\footnote{熊野大運動会という、寮のグラウンドで運動会をやる企画。夏休みって帰省している人も多いし、コンパもないし何かイベントほしいよねという思いから始まる。2021年から文化部持ち込み企画で実施されている。もう私は運営者ではないけれど、またやりたいと運営を引き継いでくれる人がいて今まで5回も開催することができています。ありがたいね。}をやったり100人でおにぎり食べたり\footnote{2021年寮祭の時計台コンパと同時開催でクスノキ前でおにぎり100個配る企画をやりました。コンパの食事提供と相性が良かったのか、これも恒例になっているみたいです。2022年から実施主体が個人でなく組織になっているので、私のときより規模が大きくなっていて大変うれしい。}。それを一緒にやろうと言ってくれる仲間に出会えました。

大学近辺にはサークルなどたくさんのコミュニティがありますが、寮ほどのフットワークの軽さがあるコミュニティはなかなか無いのではと思います。もちろん、人と人なので、相性の良い人、悪い人はいます。寮生全員と仲良くなるのは難しいと思いますが、あなたと響き合う人は絶対います。一緒に楽しいことをやりましょう。

\subsection{最後に}
寮に住んでいれば楽しいこともあればしんどいこともあります。こんな寮出てやると思ったことも一度や二度ではないです。でも、コンパで、会議で、何気ない日常で、私もあなたも寮が好きなのねと感じることができると、ここに居られてよかったな、幸せだなとなります。

一方で、熊野寮が共同生活の場であることを忘れている人を目にすると悲しい気持ちになります。例えば、みんなでやるべき当番の仕事を全くしない人、反対意見ばかりで寮を支えている人へのリスペクトが感じられない人、共有スペースを大幅に占有する人、ハラスメントじみた言動がやめられない人…。
\vskip.5\baselineskip
私・あなたの寮生活は、私、あなた、そしてほかの誰かの頑張りで成り立っています。みんなが頑張るのをやめたら寮なんて1日で腐ります。「配慮」なんていう抽象的で解釈が難しい言葉は使いたくないんだけれども、やっぱり共同生活をするうえで「配慮」って大事なんじゃないのかしらと思います。（ここでの「配慮」はモノを共有する誰かの存在を意識しリスペクトを忘れないこと、としておきます。）

親元を離れて初めての一人暮らしが熊野寮という人も多いと思いますが、あなたと一緒に住む人はあなたのお世話係ではありませんし、あなたを支配する権力者でもありません。他の多くの寮生はあなたと同じ学生で、上も下もなく対等で、ただ寮というものを通してつながり合っているだけの存在です。

楽しいばっかりの寮だと思ってたのにこんな説教じみたこと言われるなんて！とがっかりさせてしまったかもしれませんが、そういう他者と共存する大変さを乗り越えてこそ自治寮の意義があるのだと思います。熊野寮のような学生自治寮は現代社会の中ではどちらかといえば滅びゆく存在です。寮に住むこと自体が、無理だと言われがちな理想を現実にしているとも言えます。
\vskip.5\baselineskip
あなたがどのような動機をもって、熊野寮の門戸を開こうとするのかはわかりませんが、学生自治寮たるここでの生活があなたにとってすこしでも良いものでありますように。せっかくなら、いろいろな人と話して、いろいろな活動に参加してみるといいと思います！
\vskip.5\baselineskip
最後までお読みいただきありがとうございました。